\scanpage{218}
\mspage{73}

une de l'autre, entre l'ensemble $\mathrm{Cl}_R(P)$ des parties $R$-closes de $P$, et l'ensemble $\mathrm{Cl}_{\unsure{{}^{s}R|{}^{t}R}}(Q)$ des parties $\,{}^{s}R$-closes (nous dirons $R$-closes par abus de langage) de $Q$. Ces bijections sont \inserted{décroissantes, donc} des anti-\deleted{\unsure{in|t}}isomorphismes pour les relations d'ordre naturelles (d'inclusion).

\margin{\rule{0.4pt}{4\baselineskip}}%
\uline{Corollaire} 5.2. Soit $(A_i)_i$ une famille de parties \uline{$R$-closes} de $P$. Alors $\bigcap A_i$ est une partie $R$-close de $P$, et

\[
 \left(\bigcap A_i\right)^{\unsure{\cup|\smile}}=
 \underbrace{\left(\bigcup_i A_i^\cup\right)^\flat}_{
 \substack{\text{\deleted{plus petite} Sup des}\\\text{parties }R\text{-closes}\\A_i^\cup\text{ de }Q}}.
 \tag*{(5.2.1)}
\]

Mais la relation $R$ \unsure{se} ``prolonge'' en une relation $\overline R$ entre

\[
 \overline P=\unsure{\mathfrak P|\mathcal P}(P),\qquad\overline Q=\mathfrak P(Q)
 \tag*{(5.13)}
\]

par la prescription

\[
\begin{aligned}
 \deleted{(A,B)\in}\underset{\substack{\cap\\\overline P\times\overline Q}}{\overline R}
  =\big\{(A,B)\in\overline P\times\overline Q\mid{}&
  B\subset A^\cup\quad\text{i.e. }A\subset B^{\unsure{\cap|\frown}}\\
 &\text{i.e. }A\times B\subset R\big\}
\end{aligned}
\tag*{(5.14)}
\]

\deleted{\unsure{Cette}} Les ensembles $\overline P,\overline Q$ sont stables par Sup quelconques (\unsure{réunions|réunion} de parties de $P$, resp $Q$), et la relation $\overline R$ sur ces ensembles ordonnés satisfait à la condition $(D_0)$
% archive sheet number: 218
